% !TEX root = ../notes.tex

\documentclass[../notes.tex]{subfiles}

\begin{document}

\section{March 16}
Here we go.

\subsection{Weyl Conjugacy Classes}
Our next goal is to redefine our strata in terms of conjugacy classes of the Weyl group $W$ of the group $G$.
\begin{definition}
	Fix $w\in W$ and a conjugacy class $\gamma$ of a reductive group $G_p$ in characteristic $p$. Then we define
	\[\mc B_{\gamma,w}\coloneqq\left\{(g,B):\gamma\times\mc B_p:\left(B,gBg^{-1}\right)\text{ are in relative position }w\right\}.\]
	Thus, this is a subset of the flag variety $\mc B$.
\end{definition}
\begin{remark}
	One can similarly fix $g$ and only look at the Borel subgroups in $\mc B_{\gamma,w}$, which are refinements of the Springer fibers.
\end{remark}
\begin{theorem}[Lusztig~2011] \label{thm:weyl-to-unip}
	Fix $w\in W$ and a conjugacy class $\gamma$ of a reductive group $G_p$ in characteristic $p$. Suppose that $w$ has minimal length its in conjugacy class. Then there is a unique unipotent conjugacy class $\gamma$ in $\op{Conj}(U_p)$ so that $\mc B_{\gamma,w}$ is nonempty and has minimal dimension. In fact, $\gamma$ depends only on the conjugacy class of $w$.
\end{theorem}
This was settled by a long analysis by computer.
\begin{remark}
	Professor Zhiwei Yun has another construction of this map in characteristic zero.
\end{remark}
\Cref{thm:weyl-to-unip} produces the following definition.
\begin{definition}
	Fix a reductive group $G_p$ in characteristic $p$. Then we define a map $\varphi_p\colon\op{Conj}W\to\op{Conj}(U_p)$ by sending the conjugacy class of $w$ (of minimal length in its conjugacy class) to the class $\gamma$ for which $\mc B_{\gamma,w}$ is nonempty of minimal dimension.
\end{definition}
We now extend this to $\op{Conj}(U_*)$.
\begin{definition}
	For a split reductive group $G$ defined over $\ZZ$, we extend the $\varphi_p$s to a map $\varphi_*\colon\op{Conj}W\to\op{Conj}U_*$ as follows.
	\begin{itemize}
		\item If $\varphi_p(w)$ is independent of $p$, then we define $\varphi_*(w)$ as the value at any $p$.
		\item If $\varphi_p(w)$ is not independent of $p$, then it turns out to only be different at a single prime $p_0$, and we define $\varphi_*(w)=\varphi_{p_0}(w)$.
	\end{itemize}
	The fibers in $W$ are called \textit{strata}.
\end{definition}
\begin{example}
	Composing with the Springer correspondence, we see that there is thus a surjective map $\op{Conj}W\to\op{Irr}_*W$. Here are the fibers of this map for $E_8$; throughout, $A_B$ denotes the representation of dimension $A$ which appears for the first time in the $B$th symmetric power of the reflection representation of the Weyl group.
	\begin{itemize}
		\item The representation $15_{56}$ has five elements in its fiber. This kind of looks like an $A_5$.
		\item The representation $700_{28}$ has four elements in its fiber. It kind of looks like an $A_4$.
		\item The representations $1400_{29}$ or $2016_{19}$ have fibers of size four, which look like $A_2\times A_2$.
		\item There are also five representations with fiber of size three, fourteen representations with fiber of size two, and fifty-two representations with fiber of size one.
	\end{itemize}
	One makes the fibers into the graphs described above by finding paths from the most elliptic to the least elliptic conjugacy classes; it turns out that each class contains a unique most elliptic class.
\end{example}
\begin{remark}
	In fact, in characteristic two, one finds that the reductive part of the centralizer of an elliptic element in $W$ is some symplectic group, and it kind of explains how one finds the sizes of the fibers. This feature works in general, not just for $E_8$.
\end{remark}
The moral is that we can parameterize conjugacy classes by representations, which we understand better.
\begin{example}
	Let's describe the map $\varphi_*$ in type $C$. This type is classical, so $\op{Conj}(U_*)=\op{Conj}(U_2)$, and it turns out that $\varphi_*=\varphi_2$. The classes admit a combinatorial description as partitions $\nu\vdash2n$ for which an odd part appears an even number of times, and any even part which appears an even number of times has an associated number $\varepsilon\in\{0,1\}$. (The invariant $\varepsilon$ disappear in characteristic zero.) The Weyl group $W$ can be described as the permutations in $S_{2n}$ which commute with the involution $k\mapsto(2n+1-k)$. It turns out that the map $\op{Conj}(W)\to\op{Conj}(U_2)$ sends $w\in W$ to the partition given by the lengths of its cycles; for cycles of even length, it may not commute with the involution, in which case we give it $\varepsilon=1$ (and we give it $\varepsilon=0$ otherwise). 
\end{example}
\begin{example}
	In type $C$, the sequences $(d_i)$ in $\op{Sym}''C_n$ with
	\[d_0\le d_1+2\le d_2+4\le\cdots\]
	parameterize $\op{Irr}_*W$.
\end{example}


\end{document}